
%%%%%%%%%%%%%%%%%%%%%%%%%%%%%%%%%%%%%%%%
%           main body
%%%%%%%%%%%%%%%%%%%%%%%%%%%%%%%%%%%%%%%%

%-----------------------------------
%	TITLE PAGE
%-----------------------------------

\title[第四章\\
勒贝格积分]{\texorpdfstring{\LARGE \bfseries 第四章\quad 勒贝格积分}{第四章 勒贝格积分}} % The short title appears at the bottom of every slide, the full title is only on the title page
\author[\S 4.3\\
积分序列的极限] % Your institution as it will appear on the bottom of every slide, maybe shorthand to save space
{\Large \bfseries \S 4.3 积分序列的极限}
% \author{Singular Value Decomposition} % Your name
% \date{\today} % Date can be changed to a custom date
\date{}

%------------------------------------------------

\begin{document}

%------------------------------------------------

\begin{frame}

\titlepage

\end{frame}

%------------------------------------------------

\section[引言]{引言: 换序定理的谱系}

%------------------------------------------------

\begin{frame}{引言: 积分与极限交换次序}

本节研究 $\displaystyle \lim_n \int_E f_n ~ \mathrm{d} m \overset{?}{=} \int_E \lim_n f_n ~ \mathrm{d} m$ ---
勒贝格积分理论的\alert{核心地位}所在:

\begin{center}
\begin{tabular}{ccc}
\textbf{定理} & \textbf{条件} & \textbf{状态} \\[0.2em]
Levi 定理 & 渐升非负, 点态收敛 & \S 4.1 已提前引入并使用 \\
逐项积分 & 非负级数 & 本节 (Levi 的级数形式) \\
Fatou 引理 & 非负列, 无单调性 & 本节 (弱形式, 与 Levi 等价) \\
\alert{Lebesgue 控制收敛} & $|f_n| \leqslant g \in L_E$ & 本节 (最重要) \\
Lebesgue--Vitali & 依测度收敛 $+$ 等度绝对连续 & 本节 (深化) \\
\end{tabular}
\end{center}

\pause

\begin{itemize}
\item 对比 Riemann 积分: 交换次序通常要求\alert{一致收敛} ---
  勒贝格积分只要求``被控制''这类点态条件, 优势尽显;
\item 两大引擎不变: \S 4.1 Levi 定理 $+$ \S 4.2 (绝对连续性、线性性、唯一性).
\end{itemize}

\end{frame}

%------------------------------------------------

\section[逐项积分]{逐项积分定理}

%------------------------------------------------

\begin{frame}{逐项积分定理 (非负级数)}

\begin{thm}[逐项积分]
设 $\{u_n\}$ 是可测集 $E$ 上\alert{非负}可测函数列, $f = \sum_{n=1}^{\infty} u_n$, 那么
\[
\int_E f ~ \mathrm{d} m = \sum_{n=1}^{\infty} \int_E u_n ~ \mathrm{d} m .
\]
\end{thm}

\startproof[bg=yellow, fg=red](证明)
令 $f_k(x) = \sum_{n=1}^{k} u_n(x)$, 则 $\{f_k\}_{k \in \mathbb{N}}$ 构成 $E$ 上\alert{渐升}非负可测函数列, 且
\[
\lim_{k \to \infty} f_k(x) = \lim_{k} \sum_{n=1}^{k} u_n(x) = \sum_{n=1}^{\infty} u_n(x) = f(x) .
\]
于是由 Levi 定理与积分的有限线性:
\[
\int_E f ~ \mathrm{d} m
= \int_E \lim_k f_k ~ \mathrm{d} m
= \lim_k \int_E f_k ~ \mathrm{d} m
= \lim_k \sum_{n=1}^{k} \int_E u_n ~ \mathrm{d} m
= \sum_{n=1}^{\infty} \int_E u_n ~ \mathrm{d} m .
\]
\qed

\pause

\begin{itemize}
\item 级数可逐项积分, 无需一致收敛 --- ``和的可积性''与``换序''同时免费获得.
\end{itemize}

\end{frame}

%------------------------------------------------

\begin{frame}{逐项积分的推论与注记}

\begin{cor}
设 $\{E_n\}$ 是一列互不相交的可测集, $E = \bigcup_n E_n$, $f \in L_E$, 那么
\[
\int_E f ~ \mathrm{d} m = \sum_{n=1}^{\infty} \int_{E_n} f ~ \mathrm{d} m .
\]
\end{cor}

\startproof[bg=yellow, fg=red](证明)
$f = f \cdot \sum_n \chi_{E_n} = \sum_n f \chi_{E_n}$ (不交并), 由逐项积分定理
$\int_E f ~ \mathrm{d} m = \sum_n \int_E f \chi_{E_n} ~ \mathrm{d} m = \sum_n \int_{E_n} f ~ \mathrm{d} m$. \qed

\pause

\begin{itemize}
\item \S 4.2 曾用``有限可加 $+$ 绝对连续''证明过它; 现在只是逐项积分的\alert{推论};
\item 特征函数是桥梁: 积分问题 $\leftrightarrow$ 测度问题. 取 $f = 1$, 上式即\alert{测度的 $\sigma$-可加性}
  ($\int_{E_n} 1 ~ \mathrm{d} m = m E_n$); 类似的还有本章习题 5;
\graymode
\item \textbf{非负性不可少} (反例): $E = (0,1)$, $u_n(x) = n x^{n-1} - (n+1) x^n$.
  部分和 $= 1 - (k+1)x^k \to 1$, 故 $f \equiv 1$, $\int_E f ~ \mathrm{d} m = 1$;
  但 $\int_E u_n ~ \mathrm{d} m = \big[ x^n - x^{n+1} \big]_0^1 = 0$
  (Riemann 积分, $=$ Lebesgue 积分, \S 4.4), 故 $\sum_n \int_E u_n ~ \mathrm{d} m = 0 \neq 1$.
\normalmode
\end{itemize}

\end{frame}

%------------------------------------------------

\section[Fatou 引理]{Fatou 引理}

%------------------------------------------------

\begin{frame}{Fatou 引理}

\begin{thm}[Fatou 引理]
若 $\{f_n\}$ 是可测集 $E$ 上\alert{非负}可测函数列, 那么
\[
\int_E \varliminf_{n \to \infty} f_n ~ \mathrm{d} m
\leqslant
\varliminf_{n \to \infty} \int_E f_n ~ \mathrm{d} m .
\]
\end{thm}

\startproof[bg=yellow, fg=red](证明)
令 $u_n(x) = \inf_{k \geqslant n} \{ f_k(x) \}$, 则 $\{u_n\}$ 构成\alert{渐升}的非负可测函数列
(\S 3.1: $\inf$ 保可测), 且 $u_n \leqslant f_k\ (\forall k \geqslant n)$ 以及
\[
\lim_{n \to \infty} u_n(x) = \lim_{n} \inf_{k \geqslant n} \{ f_k(x) \} = \varliminf_{n \to \infty} f_n(x) .
\]
对 $\{u_n\}$ 应用 Levi 定理: $\displaystyle \int_E \varliminf_n f_n ~ \mathrm{d} m
= \lim_n \int_E u_n ~ \mathrm{d} m$; 而由 $u_n \leqslant f_k\ (k \geqslant n)$,
$\int_E u_n ~ \mathrm{d} m \leqslant \inf_{k \geqslant n} \int_E f_k ~ \mathrm{d} m$, 令 $n \to \infty$:
\[
\int_E \varliminf_{n \to \infty} f_n ~ \mathrm{d} m
= \lim_n \int_E u_n ~ \mathrm{d} m
\leqslant \lim_n \inf_{k \geqslant n} \int_E f_k ~ \mathrm{d} m
= \varliminf_{n \to \infty} \int_E f_n ~ \mathrm{d} m .
\]
\qed

\end{frame}

%------------------------------------------------

\begin{frame}{Fatou 引理的两个注记}

\begin{itemize}
\item \textbf{注 (1)}: ``$<$''可以取到. $E = [0,1]$,
\[
f_n(x) =
\begin{cases}
n, & x \in (0, \tfrac{1}{n}), \\
0, & x \in E \setminus (0, \tfrac{1}{n}) ,
\end{cases}
\qquad
\varliminf_n f_n(x) = 0 ,\ \forall x \in E,
\]
则 $\int_E \varliminf_n f_n ~ \mathrm{d} m = 0 < 1 = \lim_n \int_E f_n ~ \mathrm{d} m$
($\int_E f_n ~ \mathrm{d} m = n \cdot \tfrac1n = 1$).

\pause

\item \textbf{注 (2)}: 证明用了 Levi, 实际上二者\alert{等价} ---
Fatou $\Rightarrow$ Levi: 设 $0 \leqslant f_1 \leqslant f_2 \leqslant \cdots$,
$f = \lim_n f_n = \varliminf_n f_n$. 由 Fatou 引理
\[
\int_E f ~ \mathrm{d} m
= \int_E \varliminf_n f_n ~ \mathrm{d} m
\leqslant \varliminf_n \int_E f_n ~ \mathrm{d} m
= \lim_n \int_E f_n ~ \mathrm{d} m ;
\qquad \text{……\textcircled{1}}
\]
另一方面显然 $f \geqslant f_n \geqslant 0$, 由保序性 $\int_E f ~ \mathrm{d} m \geqslant \int_E f_n ~ \mathrm{d} m$,
对 $n$ 取极限: $\int_E f ~ \mathrm{d} m \geqslant \lim_n \int_E f_n ~ \mathrm{d} m$. ……\textcircled{2}
\end{itemize}

\pause

\begin{attnblock}
由 \textcircled{1} \textcircled{2} 得 $\int_E f ~ \mathrm{d} m = \lim_n \int_E f_n ~ \mathrm{d} m$,
即 Levi 定理 --- 两大基本定理互为等价.
\end{attnblock}

\end{frame}

%------------------------------------------------

\section[控制收敛]{Lebesgue 控制收敛定理}

%------------------------------------------------

\begin{frame}{Lebesgue 控制收敛定理}

\begin{thm}[控制收敛]
设可测集 $E$ 上的可积函数列 $\{f_n\}_{n \in \mathbb{N}}$ 满足
\begin{enumerate}
\item $f_n(x)$ 的极限存在: $f(x) = \lim_{n \to \infty} f_n(x)$, a.e.\ 于 $E$;
\item $\exists\, g \in L_E$, 使得 $|f_n(x)| \leqslant g(x)$, a.e.\ 于 $E$;
\end{enumerate}
那么 $f \in L_E$, 且有
\[
\int_E f ~ \mathrm{d} m = \lim_{n \to \infty} \int_E f_n ~ \mathrm{d} m .
\]
\end{thm}

\pause

\begin{itemize}
\item 条件 (1)(2) 可减弱为 a.e.\ 成立 (零测集上修正定义, 积分不变);
\item 勒贝格控制收敛定理是\alert{最重要的}积分与极限交换次序的结果 ---
  $|f| \leqslant g$ 这一条件比``渐升非负''更\alert{广泛}也更自然 (页边预告: 还蕴含一致收敛情形).
\end{itemize}

\end{frame}

%------------------------------------------------

\begin{frame}{控制收敛定理的证明: 两次 Fatou}

\startproof[bg=yellow, fg=red](证明)
先证 $f \in L_E$: 由 (1)(2), $|f| = \lim_n |f_n| \leqslant g$ a.e., 由 \S 4.1 控制可积性
$|f| \in L_E \Rightarrow f \in L_E$.

\pause

\textbf{1\textdegree.} 由 $|f_n| \leqslant g$ 知 $g + f_n \geqslant 0$ 是非负可测函数列. 由 Fatou 引理,
\[
\varliminf_n \int_E (g + f_n) ~ \mathrm{d} m
\geqslant \int_E \varliminf_n (g + f_n) ~ \mathrm{d} m
= \int_E (g + \varliminf_n f_n) ~ \mathrm{d} m
= \int_E g ~ \mathrm{d} m + \int_E f ~ \mathrm{d} m ,
\]
左端 $= \int_E g ~ \mathrm{d} m + \varliminf_n \int_E f_n ~ \mathrm{d} m$, 故
\[
\varliminf_{n \to \infty} \int_E f_n ~ \mathrm{d} m \geqslant \int_E f ~ \mathrm{d} m .
\qquad \text{……\textcircled{1}}
\]

\pause

\textbf{2\textdegree.} 同样 $g - f_n \geqslant 0$ 也是非负可测函数列. 由 Fatou 引理,
\[
\varliminf_n \int_E (g - f_n) ~ \mathrm{d} m
\geqslant \int_E \varliminf_n (g - f_n) ~ \mathrm{d} m
= \int_E (g - \varlimsup_n f_n) ~ \mathrm{d} m
= \int_E g ~ \mathrm{d} m - \int_E f ~ \mathrm{d} m ,
\]
左端 $= \int_E g ~ \mathrm{d} m - \varlimsup_n \int_E f_n ~ \mathrm{d} m$, 故
\[
\varlimsup_{n \to \infty} \int_E f_n ~ \mathrm{d} m \leqslant \int_E f ~ \mathrm{d} m .
\qquad \text{……\textcircled{2}}
\]

\pause

由 \textcircled{1} \textcircled{2} 以及 $\varliminf \leqslant \varlimsup$:
\[
\int_E f ~ \mathrm{d} m
= \varliminf_n \int_E f_n ~ \mathrm{d} m
= \varlimsup_n \int_E f_n ~ \mathrm{d} m
= \lim_{n \to \infty} \int_E f_n ~ \mathrm{d} m .
\]
\qed

\end{frame}

%------------------------------------------------

\begin{frame}{推论: 有界收敛定理}

\begin{cor}[有界收敛]
设 $mE < \infty$, $E$ 上可测函数列 $\{f_n\}$ 满足 $|f_n(x)| \leqslant M$
($\forall n \in \mathbb{N}, x \in E$), 且 $f(x) = \lim_n f_n(x)$, 那么 $f \in L_E$, 且有
\[
\int_E f ~ \mathrm{d} m = \lim_{n \to \infty} \int_E f_n ~ \mathrm{d} m .
\]
\end{cor}

(取控制函数 $g \equiv M$: 由 $mE < \infty$ 知 $g \in L_E$.)

\pause

\begin{remark}[推广到连续参数]
设 $\{f_\alpha\}_{\alpha \in I}$ 是可测函数族 ($I$ 具有连续统的势), 满足
$|f_\alpha(x)| \leqslant g(x) \in L_E$, 设 $\alpha_0$ 是 $I$ 的一个聚点且
$\lim_{\alpha \to \alpha_0} f_\alpha(x) = f(x)$, 那么 $f \in L_E$ 且
$\int_E f ~ \mathrm{d} m = \lim_{\alpha \to \alpha_0} \int_E f_\alpha ~ \mathrm{d} m$.
\emph{(Sketch: $\forall\, \alpha_n \to \alpha_0$, 对 $\{f_{\alpha_n}\}_{n \in \mathbb{N}}$
应用控制收敛定理.)}
\end{remark}

\end{frame}

%------------------------------------------------

\begin{frame}{注: ``控制''不可去 --- 尖峰反例}

\begin{eg}
$E = [-1, 1]$,
\[
f_n(x) =
\begin{cases}
n - n^2 |x|, & 0 < |x| < \tfrac1n, \\
0, & \tfrac1n \leqslant |x| \leqslant 1 \ \text{或}\ x = 0 ,
\end{cases}
\qquad
\int_E f_n ~ \mathrm{d} m = \frac12 \cdot \frac{2}{n} \cdot n = 1 .
\]
\begin{center}
\begin{tikzpicture}[scale=0.62]
\draw[->] (-1.3,0) -- (1.5,0) node[right] {$x$};
\draw[->] (0,0) -- (0,2.6) node[above] {$y$};
\draw[dashed, domain=0.24:1.25, smooth] plot (\x, {0.25/\x});
\draw[dashed, domain=-1.25:-0.24, smooth] plot (\x, {-0.25/\x});
\draw[thick, themecolor] (-1,0) -- (0,1) -- (1,0);
\draw[thick, themecolor] (-0.5,0) -- (0,2) -- (0.5,0);
\node[right] at (0.72,0.62) {$|y| = \tfrac{1}{4|x|}$};
\node[below] at (1.05,0) {\small $1$};
\node[above right] at (0.03,1.95) {\small $f_2$};
\node[above right] at (0.55,0.85) {\small $f_1$};
\end{tikzpicture}
\end{center}
\end{eg}

\pause

\begin{itemize}
\item 处处 $f_n(x) \to 0$ ($|x| > 0$ 时 $n > 1/|x|$ 后恒为 $0$; $x = 0$ 恒为 $0$),
  但 $\lim_n \int_E f_n ~ \mathrm{d} m = 1 \neq 0 = \int_E \lim_n f_n ~ \mathrm{d} m$;
\item 无控制的原因: $f_n(x) = -|x| \big( n - \tfrac{1}{2|x|} \big)^2 + \tfrac{1}{4|x|}$,
  故 $0 < |x| \leqslant \tfrac14$ 时 $\sup_n f_n(x) \geqslant \tfrac{1}{4|x|} - |x| \geqslant \tfrac{1}{8|x|}$,
  而 $\int_{-1/4}^{1/4} \tfrac{1}{8|x|} ~ \mathrm{d} x = \infty$ --- 任何控制函数不可积.
\end{itemize}

\end{frame}

%------------------------------------------------

\section[应用]{控制收敛定理的应用}

%------------------------------------------------

\begin{frame}{应用一: 极限计算}

\begin{eg}
$E = (0, +\infty)$, 计算
\[
\lim_{n \to \infty} \int_E \Big( 1 + \frac{x}{n} \Big)^n e^{-2x} ~ \mathrm{d} m .
\]
\end{eg}

\startproof[bg=yellow, fg=red](解)
令 $f_n(x) = \big( 1 + \tfrac{x}{n} \big)^n e^{-2x}$. 由 $\big( 1 + \tfrac{x}{n} \big)^n \leqslant e^x$,
\[
|f_n(x)| \leqslant e^x \cdot e^{-2x} = e^{-x} =: g(x) \in L_E ,
\]
且 $f_n(x) \to e^x \cdot e^{-2x} = e^{-x}$ 点态. 由控制收敛定理,
\[
\lim_{n} \int_E \Big( 1 + \frac{x}{n} \Big)^n e^{-2x} ~ \mathrm{d} m
= \int_E \lim_n \Big( 1 + \frac{x}{n} \Big)^n e^{-2x} ~ \mathrm{d} m
= \int_E e^{-x} ~ \mathrm{d} x = \big[ -e^{-x} \big]_0^{+\infty} = 1 .
\]
\qed

\pause

\begin{itemize}
\item 套路: 找\alert{可积控制} $g$ + 算\alert{点态极限} ---
  极限与积分换序的``两步走'', 是控制收敛定理的标准用法.
\end{itemize}

\end{frame}

%------------------------------------------------

\begin{frame}{应用二: Lebesgue 积分和}

Riemann 积分: \alert{定义域}分割 $a = x_0 < x_1 < \cdots < x_{n(P)} = b$, 标志点 $\xi_k$,
$\lambda(P) = \max_k \Delta x_k \to 0$,
$\int_a^b f ~ \mathrm{d}x = \lim_{\lambda(P) \to 0} \sum_{k=1}^{n(P)} f(\xi_k) \Delta x_k$.

\pause

\begin{thm}[积分和]
设 $f$ 是有界可测集 $E$ 上\alert{非负}可积函数. 任给\alert{值域}的划分
\[
P_n:\ 0 = y_0^{(n)} < y_1^{(n)} < \cdots < y_{b_n}^{(n)} < \infty,
\quad y_{b_n}^{(n)} \to \infty,\quad
\lambda_n := \max_i (y_i^{(n)} - y_{i-1}^{(n)}) \to 0 ,
\]
取水平集 $E_i^{(n)} = E\big( y_{i-1}^{(n)} \leqslant f < y_i^{(n)} \big)$ 与标志点
$\xi_i^{(n)} \in E_i^{(n)}$, 令 $\varphi_n = \sum_{i=1}^{b_n} f(\xi_i^{(n)}) \chi_{E_i^{(n)}}$, 则
\[
\int_E f ~ \mathrm{d} m = \lim_{n \to \infty} \sum_{i=1}^{b_n} f\big( \xi_i^{(n)} \big)\, m E_i^{(n)} .
\]
\end{thm}

\pause

\begin{itemize}
\item 区别在\alert{划分值域} --- 这正是第二章引入 Lebesgue 测度的动机的再现;
\item 证: $x \in E_i^{(n)}$ 时 $|\varphi_n(x) - f(x)| < \lambda_n$, 故 $\varphi_n \to f$ 点态;
  可设 $\lambda_n < 1$, 则 $\varphi_n < f + 1 \in L_E$, 由控制收敛定理即得.
\end{itemize}

\end{frame}

%------------------------------------------------

\begin{frame}{应用三: 积分号下求导}

\begin{thm}[积分号下求导]
设 $f(x, t)$ 定义在 $E \times (\alpha, \beta)$ 上, $\forall\, t \in (\alpha, \beta)$,
$f(x, t) \in L_E$ (关于 $x$); $\forall\, x \in E$, $f(x, t)$ 关于 $t$ 处处可微.
若存在 $F(x) \in L_E$, 使得
\[
\Big| \frac{\partial}{\partial t} f(x, t) \Big| \leqslant F(x),
\qquad \forall\, (x, t) \in E \times (\alpha, \beta) ,
\]
那么有
\[
\frac{d}{dt} \int_E f(x, t) ~ \mathrm{d}x
= \int_E \frac{\partial}{\partial t} f(x, t) ~ \mathrm{d}x .
\]
\end{thm}

\startproof[bg=yellow, fg=red](证明)
由积分的线性,
\[
\frac{d}{dt} \int_E f(x, t) ~ \mathrm{d}x
= \lim_{h \to 0} \int_E \frac{f(x, t+h) - f(x, t)}{h} ~ \mathrm{d}x
\overset{\text{中值定理}}{=}
\lim_{h \to 0} \int_E \frac{\partial}{\partial t} f\big( x, t + \theta(h) h \big) ~ \mathrm{d}x
\quad (0 \leqslant \theta(h) \leqslant 1) .
\]

\pause

令 $g_h(x) = \frac{\partial}{\partial t} f(x, t + \theta(h) h)$: 它可测 (差商的极限), 且
$|g_h(x)| \leqslant F(x)$, $g_h(x) \to \frac{\partial}{\partial t} f(x, t)$ ($h \to 0$).
由连续参数型控制收敛定理:
\[
\frac{d}{dt} \int_E f(x, t) ~ \mathrm{d}x
= \lim_{h \to 0} \int_E g_h(x) ~ \mathrm{d}x
= \int_E \lim_{h \to 0} g_h(x) ~ \mathrm{d}x
= \int_E \frac{\partial}{\partial t} f(x, t) ~ \mathrm{d}x .
\]
\qed

\end{frame}

%------------------------------------------------

\section[绝对收敛]{级数的逐项积分: 一般可积函数列}

%------------------------------------------------

\begin{frame}{级数的逐项积分 (一般情形)}

\begin{thm}[逐项积分]
设 $\{f_n\}$ 为 $E$ 上可积函数列, 若
$\displaystyle \sum_{n=1}^{\infty} \int_E |f_n| ~ \mathrm{d} m < \infty$
(\alert{绝对收敛}条件), 则 $\sum_{n=1}^{\infty} f_n(x)$ 在 $E$ 上几乎处处收敛. 记
$f(x) = \sum_{n=1}^{\infty} f_n(x)$, 那么 $f \in L_E$, 且有
\[
\int_E f ~ \mathrm{d} m = \sum_{n=1}^{\infty} \int_E f_n ~ \mathrm{d} m .
\]
\end{thm}

\pause

\begin{remark}
之前的逐项积分定理是对\alert{非负可测}函数列的; 这里是对\alert{一般可积}函数列的,
但要附加绝对收敛条件.
\end{remark}

\end{frame}

%------------------------------------------------

\begin{frame}{一般逐项积分的证明}

\startproof[bg=yellow, fg=red](证明)
令 $F(x) = \sum_{n=1}^{\infty} |f_n(x)|$ (非负可测). 由非负函数列的逐项积分定理,
\[
\int_E F ~ \mathrm{d} m = \sum_{n=1}^{\infty} \int_E |f_n| ~ \mathrm{d} m < \infty,
\]
即 $F \in L_E$, 故 $F(x)$ 在 $E$ 上几乎处处有限
$\Longrightarrow \sum_n f_n(x)$ 几乎处处\alert{绝对收敛}, 记 $f(x) = \sum_n f_n(x)$ (a.e.),
且 $|f(x)| \leqslant F(x)$ a.e., 故 $f \in L_E$.

\pause

令 $g_k(x) = \sum_{n=1}^{k} f_n(x)$ 为部分和函数列, 则 $g_k \to f$ a.e., 且
\[
|g_k(x)| = \Big| \sum_{n=1}^{k} f_n(x) \Big|
\leqslant \sum_{n=1}^{\infty} |f_n(x)| = F(x) .
\]
由\alert{控制收敛定理},
\[
\int_E f ~ \mathrm{d} m
= \int_E \lim_k g_k ~ \mathrm{d} m
= \lim_k \int_E g_k ~ \mathrm{d} m
= \lim_k \sum_{n=1}^{k} \int_E f_n ~ \mathrm{d} m
= \sum_{n=1}^{\infty} \int_E f_n ~ \mathrm{d} m .
\]
\qed

\end{frame}

%------------------------------------------------

\section[等度积分]{等度绝对连续积分}

%------------------------------------------------

\begin{frame}{本节最后: 更深刻的收敛定理}

本节最后, 我们建立关于积分列\alert{更深刻}的收敛定理 (去掉``控制函数''). 为此, 把 \S 4.2
单个可积函数的绝对连续性``一致化''到整个函数族:

\begin{Def}[等度绝对连续积分]
设 ($mE < \infty$) $E$ 可测, $\{f_\alpha\}_{\alpha \in I}$ 为 $E$ 上可积函数族.
若 $\forall\, \varepsilon > 0$, $\exists\, \delta > 0$, 使得 $\forall\, A \subset E$,
当 $mA < \delta$ 时不等式
\[
\int_A |f_\alpha| ~ \mathrm{d} m < \varepsilon
\]
关于 $\alpha$ \alert{一致成立} (即 $\delta$ 的选取只与 $\varepsilon$ 有关, 与 $\alpha$ 无关),
则称函数族 $\{f_\alpha\}_{\alpha \in I}$ 有\alert{等度的绝对连续积分}.
\end{Def}

\pause

\begin{itemize}
\item 单个函数: \S 4.2 绝对连续性 ($\delta$ 可依赖该函数); 族: $\delta$ 必须统一 ---
  ``等度''= uniformly;
\item 单调列 $\subset$ 控制列 $\subset$ 等度绝对连续列 --- 逐步放宽对``列''的要求.
\end{itemize}

\end{frame}

%------------------------------------------------

\begin{frame}{引理: 依测度收敛 $+$ 等度绝对连续 $\Rightarrow$ 可积}

\begin{lem}
设 $mE < \infty$, $\{f_n\}$ 是 $E$ 上可积函数列, 且\alert{依测度}收敛于 $f$, 那么当
$\{f_n\}$ 在 $E$ 上有等度的绝对连续积分时, $f \in L_E$.
\end{lem}

\begin{remark}
依测度收敛\alert{弱于}几乎处处收敛, 故加上了``一致性''条件 (等度绝对连续积分),
但也\alert{去掉了控制函数}这一条件; 结论也减弱为 $f \in L_E$ --- 没有积分与极限交换次序的结论.
\end{remark}

\startproof[bg=yellow, fg=red](证明)
由等度绝对连续性, $\forall\, \varepsilon > 0$, $\exists\, \delta(\varepsilon) > 0$,
$\forall\, A \subset E$, $mA < \delta(\varepsilon)$ 时 $\int_A |f_n| ~ \mathrm{d} m < \varepsilon$
对\alert{所有} $n$ 成立.

\pause

由 Riesz 定理 (\S 3.2), 依测度收敛列存在子列 $\{f_{n_k}\}$ a.e.\ 收敛于 $f$.
$\{f_{n_k}\}$ 仍依测度收敛, 故为依测度基本列:
$\lim_{i, j \to \infty} mE\big( |f_{n_i} - f_{n_j}| \geqslant \varepsilon \big) = 0$.
于是存在 $N(\varepsilon)$, 使得 $i, j > N(\varepsilon)$ 时,
$A = E\big( |f_{n_i} - f_{n_j}| \geqslant \varepsilon \big)$ 满足 $mA < \delta(\varepsilon)$. 按
$E(|f_{n_i} - f_{n_j}| < \varepsilon)$ 与 $A$ 分割:
\[
\int_E |f_{n_i} - f_{n_j}| ~ \mathrm{d} m
\leqslant \int_E \varepsilon ~ \mathrm{d} m
+ \int_A |f_{n_i}| ~ \mathrm{d} m
+ \int_A |f_{n_j}| ~ \mathrm{d} m
\leqslant \varepsilon \cdot mE + \varepsilon + \varepsilon
= (mE + 2)\, \varepsilon .
\]
\qed

\end{frame}

%------------------------------------------------

\begin{frame}{引理证明 (续): Fatou 收尾}

\startproof[bg=yellow, fg=red](证明续)
故 $\big\{ \int_E |f_{n_k}| ~ \mathrm{d} m \big\}_{k \in \mathbb{N}}$ 是 $\mathbb{R}$ 中的
\alert{基本列} (Cauchy 列), 极限存在有限. 对 $\{|f_{n_k}|\}$ 应用 Fatou 引理:
\[
\int_E |f| ~ \mathrm{d} m
= \int_E \varliminf_k |f_{n_k}| ~ \mathrm{d} m
\leqslant \varliminf_k \int_E |f_{n_k}| ~ \mathrm{d} m
< \infty
\ \Longrightarrow\ |f| \in L_E \Longrightarrow f \in L_E .
\]
\qed

\end{frame}

%------------------------------------------------

\section[Vitali]{Lebesgue--Vitali 定理}

%------------------------------------------------

\begin{frame}{Lebesgue--Vitali 定理}

\begin{thm}[Lebesgue--Vitali]
设 $mE < \infty$, $E$ 上可积函数列 $\{f_n\}$ \alert{依测度}收敛于 $f$, 那么
\[
\lim_{n \to \infty} \int_E |f_n - f| ~ \mathrm{d} m = 0
\quad \Longleftrightarrow \quad
\{f_n\}\ \text{在}\ E\ \text{上有等度的绝对连续积分} .
\]
\end{thm}

\pause

\startproof[bg=yellow, fg=red](证明 ``$\Leftarrow$'')
由引理, $f \in L_E$. 由依测度收敛, $\forall\, \varepsilon > 0$, $\exists\, N(\varepsilon)$,
$\forall\, n > N(\varepsilon)$:
$A = E(|f_n - f| \geqslant \varepsilon)$ 满足 $mA < \delta(\varepsilon)$ (等度性). 于是
\[
\int_E |f_n - f| ~ \mathrm{d} m
= \int_{E(|f_n - f| < \varepsilon)} |f_n - f| ~ \mathrm{d} m
+ \int_A |f_n - f| ~ \mathrm{d} m
\leqslant \varepsilon \cdot mE + \varepsilon + \varepsilon
= (mE + 2)\, \varepsilon ,
\]
(末项 $\int_A |f| ~ \mathrm{d} m \leqslant \varepsilon$ 用 $f \in L_E$ 的绝对连续性).
故 $\lim_n \int_E |f_n - f| ~ \mathrm{d} m = 0$.

\end{frame}

%------------------------------------------------

\begin{frame}{Lebesgue--Vitali 定理证明 (续)}

\startproof[bg=yellow, fg=red](证明 ``$\Rightarrow$'' 续)
设 $\lim_n \int_E |f_n - f| ~ \mathrm{d} m = 0$. 先看 $f$: 取定一个 $n$,
$\int_E |f| ~ \mathrm{d} m \leqslant \int_E |f_n - f| ~ \mathrm{d} m + \int_E |f_n| ~ \mathrm{d} m < \infty$,
即 $f \in L_E$, 从而 $|f|$ 的积分绝对连续: $\forall\, \varepsilon > 0$, $\exists\, \delta_0(\varepsilon)$,
$mA < \delta_0(\varepsilon) \Rightarrow \int_A |f| ~ \mathrm{d} m < \varepsilon$.

\pause

考察 $\int_A |f_n| ~ \mathrm{d} m \leqslant \int_A |f_n - f| ~ \mathrm{d} m + \int_A |f| ~ \mathrm{d} m$:
\begin{itemize}
\item 当 $n > N(\varepsilon)$ 时
  ($\int_E |f_n - f| ~ \mathrm{d} m < \varepsilon$):
  $\int_A |f_n| ~ \mathrm{d} m \leqslant \int_E |f_n - f| ~ \mathrm{d} m + \varepsilon \leqslant 2\varepsilon$;
\pause
\item 对 $n = 1, 2, \cdots, N(\varepsilon)$ (有限个), 由 $|f_n - f| \in L_E$ 的绝对连续性, 分别存在
  $\delta_1(\varepsilon), \cdots, \delta_{N(\varepsilon)}(\varepsilon)$, 使
  $mA < \delta_n(\varepsilon) \Rightarrow \int_A |f_n - f| ~ \mathrm{d} m < \varepsilon$
  $\Rightarrow \int_A |f_n| ~ \mathrm{d} m \leqslant 2\varepsilon$.
\end{itemize}

\pause

取
\[
\delta(\varepsilon) = \min \big\{ \delta_0(\varepsilon), \delta_1(\varepsilon), \cdots,
\delta_{N(\varepsilon)}(\varepsilon) \big\} ,
\]
则 $\forall\, A \subset E$, $mA < \delta(\varepsilon)$ 时, $\int_A |f_n| ~ \mathrm{d} m \leqslant 2\varepsilon$
对\alert{一切} $n$ 成立, 即 $\{f_n\}_{n \in \mathbb{N}}$ 具有等度的绝对连续积分. \qed

\pause

\begin{attnblock}
注: ``$\Rightarrow$''部分不要求 $mE < \infty$; 但 ``$\Leftarrow$''部分以及之前的引理需要 $mE < \infty$,
否则有反例 (下一页).
\end{attnblock}

\end{frame}

%------------------------------------------------

\begin{frame}{反例: $mE = \infty$ 时结论失效}

\graymode

\begin{eg}
$E = (0, \infty)$, $f_n = \dfrac{1}{n} \chi_{(0, n)}$, $f = 0$.
\end{eg}

\begin{enumerate}
\item \textbf{依测度收敛}: $\forall\, \varepsilon > 0$, 取 $N = \lfloor \tfrac1\varepsilon \rfloor + 1$,
  则 $n > N$ 时 $f_n(x) \leqslant \tfrac1n < \varepsilon$ 处处成立, 即
  $E(|f_n - f| \geqslant \varepsilon) = \emptyset$,
  $\lim_n mE(|f_n - f| \geqslant \varepsilon) = 0$;
\item \textbf{等度绝对连续}: $|f_n| \leqslant 1$ 且 $\forall\, \varepsilon > 0$ 取 $\delta = \varepsilon$:
  $mA < \delta$ 时 $\int_A |f_n| ~ \mathrm{d} m \leqslant \int_A ~ \mathrm{d} m = mA < \varepsilon$;
\item \textbf{但}: $\int_E f ~ \mathrm{d} m = 0$,
  而 $\int_E f_n ~ \mathrm{d} m = \tfrac1n \cdot n = 1 \not\to 0$ ---
  ``$\Leftarrow$''在 $mE = \infty$ 时失效.
\end{enumerate}

\pause

\begin{itemize}
\item 引理的反例: 取 $f_n = \sum_{k=0}^{n} \frac{1}{k+1} \chi_{[k, k+1)}$
  (等度绝对连续 $+$ 依测度收敛, 但极限函数 $f = \sum_{k \geqslant 0} \frac{\chi_{[k,k+1)}}{k+1} \notin L_E$).
\item 批注: 课本上相应例题不清晰, 这里采用上述构造.
\end{itemize}

\normalmode

\end{frame}

%------------------------------------------------

\begin{frame}{推论: 依测度收敛型控制收敛定理}

\begin{cor}[依测度型控制收敛]
设可测集 $E$ 上可测函数列 $\{f_n\}_{n \in \mathbb{N}}$ 满足
\begin{enumerate}
\item $\{f_n\}_{n \in \mathbb{N}}$ \alert{依测度}收敛于可测函数 $f$;
\item $\exists\, g \in L_E$, 使得 $|f_n(x)| \leqslant g(x)$, a.e.\ $x \in E$;
\end{enumerate}
那么 $f \in L_E$, 且有 $\displaystyle \int_E f ~ \mathrm{d} m = \lim_{n \to \infty} \int_E f_n ~ \mathrm{d} m$.
\end{cor}

\startproof[bg=yellow, fg=red](证明 1\textdegree: $mE < \infty$)
由 $|f_n| \leqslant g$ 及 $g$ 的积分绝对连续性: $\exists\, \delta(\varepsilon)$,
$mA < \delta$ 时 $\int_A |f_n| ~ \mathrm{d} m \leqslant \int_A g ~ \mathrm{d} m < \varepsilon$
($\forall n$) --- $\{f_n\}$ 有等度绝对连续积分. 由 Lebesgue--Vitali 定理:
\[
\lim_n \int_E |f_n - f| ~ \mathrm{d} m = 0 .
\qquad \text{……(*)}
\]
由 $(*)$: $\int_E |f| ~ \mathrm{d} m \leqslant \int_E |f_n - f| ~ \mathrm{d} m + \int_E |f_n| ~ \mathrm{d} m$
取 $n$ 充分大即有限, 故 $f \in L_E$; 且
$|\int_E f ~ \mathrm{d} m - \int_E f_n ~ \mathrm{d} m| \leqslant \int_E |f_n - f| ~ \mathrm{d} m \to 0$,
即 $\int_E f ~ \mathrm{d} m = \lim_n \int_E f_n ~ \mathrm{d} m$.

\end{frame}

%------------------------------------------------

\begin{frame}{依测度型控制收敛: 无穷测度的情形}

\startproof[bg=yellow, fg=red](证明 2\textdegree 续: $mE = \infty$)
由 $g \in L_E$, $\forall\, \varepsilon > 0$, $\exists\, X \in \mathbb{R}_+$, 使得
$\alpha > X$ 时 $\int_{E \setminus (-\alpha, \alpha)} g ~ \mathrm{d} m < \varepsilon$.
于是 (截断, $\alpha$ 只与 $\varepsilon$ 有关、与 $n$ 无关):
\[
\int_E |f_n - f| ~ \mathrm{d} m
= \Big( \int_{E \setminus (-\alpha, \alpha)} + \int_{(-\alpha, \alpha) \cap E} \Big) |f_n - f| ~ \mathrm{d} m
\leqslant \int_{E \setminus (-\alpha, \alpha)} 2 g ~ \mathrm{d} m
+ \int_{(-\alpha, \alpha) \cap E} |f_n - f| ~ \mathrm{d} m ,
\]
(其中 $|f_n|, |f| < g$ a.e.) 第一项 $\leqslant 2\varepsilon$; 第二项由 1\textdegree
(在有限测度集 $(-\alpha, \alpha) \cap E$ 上) 当 $n \to \infty$ 时 $\to 0$. 故
\[
\varlimsup_n \int_E |f_n - f| ~ \mathrm{d} m \leqslant 2\varepsilon
\ \xrightarrow{\varepsilon \to 0}\
\lim_{n \to \infty} \int_E |f_n - f| ~ \mathrm{d} m = 0 ,
\]
再由前一页的论证 ($\int_E |f| \leqslant \int_E |f_n - f| + \int_E |f_n|$,
$|\int_E f - \int_E f_n| \leqslant \int_E |f_n - f|$) 得
$f \in L_E$ 且 $\int_E f ~ \mathrm{d} m = \lim_n \int_E f_n ~ \mathrm{d} m$. \qed

\pause

\begin{attnblock}
注: 实际上, (点态) 控制收敛定理的证明给出的是\alert{更强}的结论
$\lim_n \int_E |f_n - f| ~ \mathrm{d} m = 0$ ---
即 $\{f_n\}$ 在 $E$ 上\alert{依 $L^1$ 意义}收敛于 $f$ (第五章的主角).
\end{attnblock}

\end{frame}

%------------------------------------------------

\section[小结]{小结与展望}

%------------------------------------------------

\begin{frame}{小结与展望}

\begin{itemize}
\item \textbf{换序谱系}: 逐项积分 (非负级数) $\to$ Fatou ($\varliminf$ 弱形式, 与 Levi 等价)
  $\to$ \alert{控制收敛} ($|f_n| \leqslant g \in L_E$, 最重要) $\to$ 有界收敛 ($mE < \infty$);
  条件逐步放宽: 渐升 $\Rightarrow$ 控制 $\Rightarrow$ 等度绝对连续 $+$ 依测度;
\pause
\item \textbf{三个应用}: 极限计算 (找控制 $+$ 算点态极限); Lebesgue 积分和 (划分\alert{值域},
  呼应测度动机); 积分号下求导 (中值定理 $+$ 连续参数型);
\pause
\item \textbf{级数版}: $\sum_n \int_E |f_n| ~ \mathrm{d} m < \infty \Rightarrow$
  逐项积分 (绝对收敛换非负性);
\pause
\item \textbf{深化}: 等度绝对连续积分 (一致化 \S 4.2); Lebesgue--Vitali:
  依测度收敛下 $L^1$ 收敛 $\iff$ 等度绝对连续 ($mE < \infty$);
  依测度型控制收敛 ($mE = \infty$ 用截断);
\pause
\item \textbf{注}: (点态) 控制收敛实际证明了 $L^1$ 收敛 ---
  $\lim_n \int_E |f_n - f| ~ \mathrm{d} m = 0$, 为第五章 $L^p$ 空间埋下伏笔;
\pause
\item \textbf{展望}: \S 4.4 Riemann 积分与 Lebesgue 积分的比较 ---
  $L$ 积分是 $R$ 积分的推广 (Lebesgue 判别法), 布置作业: 第四章 21, 23, 25, 26.
\end{itemize}

\end{frame}

%------------------------------------------------

\end{document}
